\noindent Sabemos que el algoritmo de muestreo por rechazo es:

\medskip
\begin{ttfamily}
  \noindent\textbf{fun} SAMPLING$(X, Y, bn, N)$\\
  \textbf{INPUT:} $X$-obs, $Y$-pred, $bn$-red, $N$-num.\ de muestras\\[1mm]
  \hspace*{4mm} $C \leftarrow \text{ZEROES}(\text{size} = |Y|)$\\
  \hspace*{4mm} \textbf{for} $(j$ \textbf{from} $1$ \textbf{to} $N)$ \textbf{do:}\\
  \hspace*{8mm} $x \leftarrow \text{PRIOR\_SAMPLING}(bn)$\\
  \hspace*{8mm} \textbf{if} $(x$ es consistente con $X)$ \textbf{then:}\\
  \hspace*{12mm} $C[x] \leftarrow C[x] + 1$\\
  \hspace*{4mm} \textbf{return} $\text{NORM}(C)$
\end{ttfamily}

\bigskip
\noindent\textit{Básicamente:}
\begin{enumerate}
  \item Genera muestras $\;\to\;$ ya nos las dan
  \item Rechaza las que no coinciden con la evidencia
  \item Con las que sí coincide, cuenta cuántas veces aparece cada valor $x$ de la variable
  \item Normaliza
\end{enumerate}

\noindent Entonces, si nuestra evidencia es $X_1 = T,\; X_3 = F,\; X_4 = T$ y queremos inferir $X_2$, vemos cada muestra y rechazamos si no coinciden con la evidencia:

\medskip
\begin{itemize}
  \item[a)] $X_1{=}T,\; X_2{=}F,\; X_3{=}F,\; X_4{=}T$ \quad Coincide, no rechazamos.
  \item[b)] $X_1{=}T,\; X_2{=}T,\; X_3{=}F,\; X_4{=}T$ \quad Coincide, no rechazamos.
  \item[c)] $X_1{=}T,\; X_2{=}F,\; X_3{=}F,\; X_4{=}T$ \quad Coincide, no rechazamos.
  \item[d)] $X_1{=}T,\; X_2{=}F,\; X_3{=}F,\; X_4{=}T$ \quad Coincide, no rechazamos.
  \item[e)] $X_1{=}F,\; X_2{=}F,\; X_3{=}F,\; X_4{=}T$ \quad No coincide. $X_1 = T$ en evidencia, rechazamos.
\end{itemize}

\noindent Tenemos 4 aceptadas, contamos:
\[
  \begin{aligned}
    \text{count}(X_2 = T) &= 1\\
    \text{count}(X_2 = F) &= 3
  \end{aligned}
\]

\noindent Normalizamos siguiendo:
\[
  P(Y = y_i) = \frac{C[y_i]}{\displaystyle\sum_j C[y_j]}
  \quad
  \text{donde } C[y_i] = \text{conteo de ese valor, } \textstyle\sum_j C[y_j] = \text{suma de todos los conteos}
\]

$\therefore$ La estimación de $X_2$ es:
\[
  \Rightarrow\quad
  P(X_2 = T) = \frac{1}{1+3} = \frac{1}{4} = 0.25
\]
\[
  \Rightarrow\quad
  P(X_2 = F) = \frac{3}{1+3} = \frac{3}{4} = 0.75
\]

$\hfill \blacksquare$
